\chapter{Einleitungsanalysen}
\section{Rahmenbedingungen \& Vorkenntnisse}
\subsection{Rahmenbedingungen}
Die Unterrichtseinheit richtet sich an eine GYM-3-Klasse, die das Ergänzungsfach Geografie besucht. Die Klasse umfasst ca. 20 motivierte Schüler:innen, die ein vertieftes Interesse an geographischen und wirtschaftlichen Fragestellungen mitbringen. Der Unterricht findet in einem gewöhnlichen Klassenzimmer mit allen notwendigen Ressourcen (Tafel, Beamer, Internetzugang, Karten, Lehrbuch Arbeitsblätter) statt.

\subsection{Vorkenntnisse}
Die Schüler:innen haben bereits grundlegende Kenntnisse zu:
\begin{itemize}
    \item Humangeografischen Konzepten (z. B. Mobilität, Siedlung, Wirtschaft, Entwicklung)
    \item Wirtschaftsgeografie (Standortfaktoren, Lieferketten, Welthandel)
    \item Globale Herausforderungen (Klimawandel, Migration, Nord-Süd-Gefälle)
    \item Nachhaltige Entwicklung (Dimensionen Wirtschaft, Gesellschaft und Umwelt) und der Einfluss menschlicher Aktivitäten auf die Umwelt via Klimawandel
\end{itemize}

Diese Vorkenntnisse werden nun auf das spezifische Thema \enquote{Globalisierung}, mit Fokus auf den Bereich \enquote{globale Lieferketten und nachhaltige Wirtschaftsprozesse} angewandt und vertieft.

\section{Sachanalyse}
Die Globalisierung beschreibt die zunehmende wirtschaftliche, kulturelle, gesellschaftliche und politische Verflechtung von Staaten und Regionen. Fortschritte in Transport- und Kommunikationstechnologien haben diesen Prozess erheblich beschleunigt und ermöglichen eine weltweite Vernetzung von Märkten, Unternehmen und Konsument:innen. Besonders deutlich wird dies in globalen Lieferketten, bei denen die Produktion eines einzigen Produkts über mehrere Länder verteilt ist.

\subsection{Globale Lieferketten und gegenseitige wirtschaftliche Abhängigkeiten}
Internationale Unternehmen nutzen die globale Arbeitsteilung, indem sie Produktionsschritte gezielt dorthin auslagern, wo sie am kostengünstigsten oder effizientesten durchgeführt werden können. Dies führt zu komplexen wirtschaftlichen Abhängigkeiten zwischen Produzent:nnen, Zulieferer:innen, Logistikunternehmen, Handelsfirmen und Konsument:innen. In der ersten Doppellektion analysieren die Schüler:innen solche Lieferketten und erarbeiten sich ein Verständnis für die vielschichtigen Akteur:innen, die an der Produktion eines Produkts beteiligt sind. Diese Analyse bildet die Grundlage für weiterführende Fragestellungen zu Machtverhältnissen, Nachhaltiger Entwicklung und Fairness in globalen Handelsstrukturen.

\subsection{Wirtschaftliche Akteur:innen und Mechanismen der Globalisierung}
Die wirtschaftliche Globalisierung wird nicht nur durch Unternehmen und Konsument:innen gesteuert, sondern auch durch eine Vielzahl von Akteur:innen mit unterschiedlichen Interessen. Neben der Wirtschaft (multinationale Unternehmen, Märkte, Standortförderungsorganisationen) spielen auch politische Akteur:innen (Staaten, internationale Organisationen, Handelsabkommen) sowie gesellschaftliche Faktoren (Arbeitsbedingungen, soziale Bewegungen, Umweltauflagen) eine entscheidende Rolle. Zur systematischen Analyse werden exogene und endogene Globalisierungstreiber unterschieden:
\begin{itemize}
    \item \textbf{Exogene Treiber:} Technologische oder wirtschaftliche Entwicklungen, die nicht direkt aus der Gesellschaft entstehen (z. B. Digitalisierung, Automatisierung, Innovationen im Transportwesen).
    \item \textbf{Endogene Treiber:} Entwicklungen innerhalb gesellschaftlicher und politischer Systeme (z. B. Handelsabkommen, Zölle, Arbeitsgesetze, geopolitische Machtstrukturen).
\end{itemize}
Die Auseinandersetzung mit diesen Akteur:innen und Einflussfaktoren erfolgt in der zweiten Doppellektion durch die Analyse von Global Cities, Global Players und internationalen Investitionen (ADI, Greenfield- und Brownfield-Investitionen). Die Schüler:innen erkennen, wie wirtschaftliche Macht auf globaler Ebene verteilt ist und welche Mechanismen die Globalisierung vorantreiben.

\subsection{Globalisierung und soziale Herausforderungen}
Die Schüler:innen zeigen in der Regel eine kritische Haltung gegenüber der Globalisierung, insbesondere in Bezug auf soziale Ungleichheiten und ausbeuterische Praktiken internationaler Unternehmen. Während globale Lieferketten wirtschaftliche Vorteile wie kostengünstige Produktion und weltweite Warenverfügbarkeit bieten, gehen sie oft mit sozialen und ökologischen Herausforderungen einher:
\begin{itemize}
    \item Geringe Löhne und schlechte Arbeitsbedingungen in Produktionsländern.
    \item Umweltbelastungen durch lange Transportwege und ressourcenintensive Produktionsprozesse.
    \item Verlagerung ökologischer und sozialer Probleme in Länder mit schwächeren Regulierungssystemen.
\end{itemize}

Da auch der Protektionismus zahlreiche Nachteile mit sich bringt, setzen sich die Schüler:innen in der dritten Doppellektion mit alternativen Ansätzen auseinander, um globale Lieferketten gesellschafts- und umweltverträglicher zu gestalten. Dabei werden verschiedene Konzepte einer „Global Green Economy“ diskutiert, um Maßnahmen zur nachhaltigen Gestaltung globaler Wertschöpfungsketten in den Dimensionen Gesellschaft, Wirtschaft und Umwelt zu entwickeln.

\subsection{Globalisierung vs. Protektionismus – Gewinner:innen und Verlierer:innen}
In der vierten Doppellektion wird der Protektionismus als Gegenentwurf zur Globalisierung thematisiert. Die Schüler:innen analysieren dessen kurz- und langfristige Folgen auf Mikro- und Makroebene und vergleichen protektionistische Maßnahmen mit einer globalisierungsfreundlichen Wirtschaftspolitik und erkennen, dass beide Systeme Gewinner:innen und Verlierer:innen produzieren. Durch die Auseinandersetzung mit diesen Gegenmodellen und das vertiefte Verständnis von Chancen und Risiken beider Ansätze sollen die Schüler:innen eine differenzierte Meinung zur Globalisierung und zum Protektionismus entwickeln.

Die Unterrichtseinheit verfolgt das Ziel, den Schüler:innen eine differenzierte Perspektive auf die Globalisierung zu vermitteln. Durch die eigenständige Analyse globaler Lieferketten, wirtschaftlicher Akteur:innen und politischer Gegenmodelle sollen sie erkennen, dass die Globalisierung weder ausschließlich positiv noch ausschließlich negativ ist, sondern eine Vielzahl von wirtschaftlichen, gesellschaftlichen und ökologischen Faktoren umfasst. Indem sie sich in unterschiedliche Rollen und Perspektiven begeben, verstehen die Schüler:innen, dass verschiedene Akteur:innen unterschiedlich von der Globalisierung profitieren – oder benachteiligt werden – und dass wirtschaftspolitische Maßnahmen stets Gewinner:innen und Verlierer:innen erzeugen.

\section{Didaktische Analyse}
\subsection{Gesellschaftsrelevanz}
Die Globalisierung prägt nahezu alle Lebensbereiche, von der Produktion und Verteilung von Waren bis hin zur politischen und kulturellen Auswirkungen. Die lokale, regionale und internationale Verlagerung von Produktionsstätten, globalen Handelsabkommen und transnationalen Unternehmen beeinflussen Wirtschaftsstrukturen und Arbeitsmärkte weltweit. In einer Zeit, in der Nachhaltige Entwicklung und faire Lieferketten zunehmend diskutiert werden, ist es essenziell für die Schüler:innen, die Mechanismen der Globalisierung und ihre Auswirkungen zu verstehen. Themen wie Fast Fashion, Elektroschrott, $_2$-Emissionen durch globale Transporte führen zu Debatten, in denen informierte Bürger:innen kritische Konsumentscheidungen treffen müssen.

\subsection{Schüler:innenrelevanz}
Die Globalisierung betrifft die Schüler:innen unmittelbar: ihre Kleidung, Lebensmittel, Technologie und Medienkonsum hängen von weltweiten Produktions- und Lieferprozessen ab. Durch die Beschäftigung mit globalen Lieferketten erhalten sie einen Einblick in wirtschaftliche Zusammenhänge im Welthandel und können bewusster über ethischen Konsum, Fair Trade und nachhaltige Alternativen nachdenken. Die Fragestellungen regen zur Reflexion über ihre eigene Rolle als Konsument:innen an:
\begin{quote}
„Welche sozialen und ökologischen Folgen hat mein Konsum?“\\
„Wie könnte eine nachhaltigere Wirtschaft auf unterschiedlichen räumlichen Massstabebenen (lokal, regional, (inter-)national, global) aussehen?“
\end{quote}
Durch interaktive Methoden wie Concept-Maps, Fallstudien und Debatten können Schüler und Schülerinnen ihre analytischen Fähigkeiten und Urteilsfähigkeit schulen.

\subsection{Fachrelevanz}
Die Auseinandersetzung mit Globalisierung und globalen Lieferketten ist zentral in der Wirtschaftsgeografie. Sie ermöglicht eine interdisziplinäre Betrachtung von Wirtschaft, Umwelt und Gesellschaft und bietet einen Raumbezug auf lokaler, regionaler, (inter-)nationaler und globaler Ebene. Im gymnasialen Unterricht fördert das Thema die Kompetenz, geografische Problemstellungen zu analysieren, (Einfluss-)Faktoren und Wechselwirkungen zu verstehen und nachhaltige Lösungsansätze zu diskutieren.

Konkret knüpft die Unterrichtseinheit an folgende Fachbereiche der Geografie an:
\begin{itemize}
    \item \textbf{Wirtschaftsgeografie:} Produktions- und Standortfaktoren, internationale Arbeitsteilung
    \item \textbf{Sozialgeografie:} Ungleichheiten im Welthandel, Fair Trade, Nachhaltige Entwicklung
    \item \textbf{Politische Geografie:} Handelsabkommen, Einfluss transnationaler Konzerne
    \item \textbf{Umweltgeografie:} CO$_2$-Emissionen, Ressourcenverbrauch, Recycling
\end{itemize}

\subsection{Räumliche Bedeutung}
Die Globalisierung hat lokale, regionale und globale Auswirkungen, die durch Fallstudien verdeutlicht werden können:
\begin{itemize}
    \item \textbf{Lokal:} Woher stammen die Produkte in einem Schweizer Supermarkt?
    \item \textbf{Regional:} Auswirkungen der Textilindustrie in Bangladesch oder der Elektronikproduktion in China
    \item \textbf{National:}  Die Rolle der Schweiz als Handelsplatz und Finanzzentrum, z. B. durch die Analyse von Handelsabkommen mit der EU
    \item \textbf{Inter-national:}  Handelsabkommen der Schweiz mit anderen Ländern, z. B. Schweiz-Indonesien-Abkommen
    \item \textbf{Global:} Die Rolle von Handelsabkommen, internationale Handelsströme, Umweltfolgen
\end{itemize}

Das Thema kann durch räumliche Visualisierungen (z. B. Karten zu Handelsrouten, Diagramme zur CO$_2$-Bilanz von Produkten) unterstützt werden, um die Komplexität der Globalisierung verständlich darzustellen.

\subsection{Zusammenfassung}
Die Unterrichtseinheit zur Globalisierung ist sowohl fachlich als auch gesellschaftlich relevant und bietet den Schüler:innen einen direkten Alltagsbezug. Durch eine interaktive, problemorientierte Herangehensweise werden sie dazu befähigt, kritische Fragen zur Wirtschafts- und Umweltverantwortung zu stellen. Der Fokus auf globale Lieferketten ermöglicht eine anschauliche, greifbare Auseinandersetzung mit den Mechanismen der Globalisierung und deren Auswirkungen.

\section{Lernzielanalyse}
Im Lehrplan 17 werden viele der genannten Kompetenzen explizit erwähnt, insbesondere:
\begin{itemize}
    \item „Verstehen die zunehmende Verflechtung von staatlichen, wirtschaftlichen, politischen und kulturellen Ebenen und die daraus resultierenden Veränderungen.“
    \item „Erkennen die Dynamik der Globalisierung und beurteilen Ursachen und Folgen der internationalen Vernetzung.“
    \item „Analysieren an Fallstudien Verflechtungen und Abhängigkeiten im Welthandel und deren Auswirkungen auf Umwelt, Gesellschaft und Wirtschaft.“
    \item „Reflektieren globale Verflechtungen in Politik, Wirtschaft, Kultur und Umwelt.“
\end{itemize}

\subsection{Orientierung an den Kompetenzzielen des Lehrplans 17}
Die Unterrichtseinheit zur Globalisierung und zu globalen Lieferketten orientiert sich an den übergeordneten Zielen des Lehrplans 17 für das Fach Geografie. Besonders relevant sind die folgenden Kompetenzbereiche:
\begin{itemize}
    \item „Die zunehmende Verflechtung von staatlichen, wirtschaftlichen, politischen und kulturellen Ebenen und die daraus resultierenden Veränderungen verstehen.“\\
    Die Schüler:innen erkennen, wie die Globalisierung Produktions-, Konsum- und Handelsprozesse beeinflusst.
    \item „Die Dynamik der Globalisierung erfassen und Ursachen sowie Folgen der internationalen Vernetzung beurteilen.“\\
    Durch die Analyse von Lieferketten lernen sie, wirtschaftliche Zusammenhänge kritisch zu reflektieren.
    \item „An Fallstudien Verflechtungen und Abhängigkeiten im Welthandel und deren Auswirkungen auf Umwelt, Gesellschaft und Wirtschaft analysieren.“\\
    Die Schüler:innen untersuchen reale Beispiele, um globale Wechselwirkungen zu verstehen.
    \item „Globale Verflechtungen in Politik, Wirtschaft, Kultur und Umwelt reflektieren.“\\
    Sie setzen sich mit den Vor- und Nachteilen der Globalisierung auseinander und bewerten deren gesellschaftliche Folgen.
\end{itemize}

\subsection{Lernziele}
\begin{itemize}
    \item (AFB I) Die SuS können die Grundbegriffe der Globalisierung (Welthandel, Vernetzung, Arbeitsteilung, Freihandel) erläutern und in einer Concept Map darstellen.
    \item (AFB II-III) Sie können eine globale Lieferkette am Beispiel eines Textil-, Technologie- oder Lebensmittelprodukts analysieren und Akteur:innen sowie deren Rollen identifizieren.
    \item (AFB II-III) Sie sind in der Lage, Hypothesen über Chancen und Risiken der Globalisierung für verschiedene Akteur:innen zu formulieren und zu diskutieren.
    \item (AFB II-III) Sie können die räumlichen Auswirkungen globaler Wirtschaftsprozesse auf verschiedene Regionen der Welt analysieren.
    \item (AFB II-III) Die SuS können aktuelle Gegenentwürfe zur Globalisierung, insbesondere protektionistische Massnahmen, hinsichtlich deren Folgen für die Wirtschaft, Gesellschaft und Umwelt kritisch betrachten und einschätzen.
    \item (AFB I-II) Die SuS verstehen die historische Entwicklung der Globalisierung (Kolonialismus, Industrialisierung, digitale Vernetzung) und deren Einfluss auf heutige wirtschaftliche Prozesse.
    \item (AFB I) Sie kennen die zentralen Akteur:innen der Globalisierung (multinationale Konzerne, Regierungen, Konsument:innen, NGOs) und deren Interessen.
    \item (AFB I) Sie verstehen die Mechanismen globaler Lieferketten und deren wirtschaftliche, soziale und ökologische Auswirkungen.
    \item (AFB II) Sie erkennen die Zusammenhänge zwischen Globalisierung, Handelsabkommen und regionalen Wirtschaftsstrukturen.
    \item (AFB II-III) Sie kennen Möglichkeiten, wie globale Lieferketten nachhaltiger gestaltet werden können, indem sie Strategien der Effizienz, Suffizienz und Konsistenz anwenden.
    \item (AFB II-III) Die SuS reflektieren ihre Rolle als Konsumentinnen in globalen Lieferketten und entwickeln ein Bewusstsein für nachhaltigen Konsum.
    \item (AFB II-III) Die erkennen soziale und ökologische Herausforderungen der Globalisierung (z. B. Arbeitsbedingungen, Umweltverschmutzung) und hinterfragen diese kritisch.
    \item (AFB III) Sie sind fähig, konstruktive Lösungen für eine nachhaltigere Wirtschaft zu diskutieren und ethische Aspekte in wirtschaftlichen Entscheidungsprozessen zu berücksichtigen.
    \item Sie arbeiten kooperativ in Gruppen, respektieren verschiedene Perspektiven und argumentieren sachlich in Diskussionen.
\end{itemize}

\subsection{Didaktische Umsetzung und Methodeneinsatz}
Die Lernziele werden durch eine strukturierte Kombination aus Inputs, Gruppenarbeiten, Diskussionen und Fallstudien umgesetzt. Dabei erfolgt die kognitive Aktivierung der Schüler:innen durch die eigenständige Analyse globaler Prozesse und deren Bewertung aus verschiedenen Perspektiven. Die Unterrichtseinheit ist adaptiv gestaltet, sodass unterschiedliche Lernniveaus berücksichtigt werden können. Dies geschieht durch offene Aufgabenstellungen und den Einsatz differenzierter Materialien, die individuelle Zugänge ermöglichen. Anwendungsaufgaben und Reflexionsphasen fördern zudem die Transferleistung, sodass die Schüler:innen ihre Erkenntnisse auf weitere wirtschaftliche Prozesse übertragen und kritisch hinterfragen können.